Vision-language models can describe scientific figures well while behaving unreliably when evidence is missing or misleading. We present \textsc{SciFig-Eval}, a diagnostic benchmark for scientific figure understanding that evaluates perception, reasoning, and behaviour rather than collapsing them into a single accuracy score. The benchmark contains 250 annotated arXiv figures and more than 23,000 model-output evaluations across eight models, combining checklist-based MQM scoring, 1,000 targeted capability questions, image transformations, poisoned captions, false-premise probes, and selective-blur tests. We introduce the Admittance--Resistance--Inductance (A-R-I) framework to measure whether models acknowledge unreadable evidence, resist misleading context, and infer only when partial visual information supports it. The results reveal sharp behavioural divergences among models with similar perception scores. GPT-5.2 achieves the best description quality (MQM 91.6) but fabricates answers about unreadable elements in 94\% of cases. Gemini 3.1 Pro scores similarly on perception (90.2), yet admits uncertainty in 90\% of such cases and leads on resistance (0.91). These findings show that perception and reasoning benchmarks alone can miss deployment-critical failure modes in scientific workflows.
